\scanpage{260}
\mspage{- 99 -}\hfill IX

que l'on a une \uline{immersion ouverte} canonique
\[
A^{\handwritten{\natural}}_{K,K'}\handwritten{\hookrightarrow}A^{\handwritten{\natural}}_{K,K''}\ ,\tag*{(8.3.1)}
\]
permettant d'identifier le premier groupe à un sous-groupe ouvert du deuxième. Pour $K'$ variable, on voit que les $A_{K,K'}$ forment ainsi un système inductif de groupes lisses de type fini sur $K$, à morphismes de transition des immersions ouvertes. Si on pose
\[
\Phi_{K,K'}=A^{\handwritten{\natural}}_{K,K'}/A^{\handwritten{\natural}\,\unsure{o|0|\circ}}_{K,K'}\ ,\tag*{(8.3.2)}
\]
les $\Phi_{K,K'}$ eux-mêmes forment donc un système inductif de groupes finis étales sur $K$, qui sera étudié plus en détail plus bas.

Nous poserons
\[
A^{\handwritten{\unsure{\flat|b}}}_K=\underset{K'}{\mathop{\lim}\limits_{\handwritten{\longrightarrow}}}A^{\handwritten{\natural}}_{K,K'}\ ,\tag*{(8.3.3)}
\]
la limite étant prise dans la catégorie des faisceaux fppf sur $K$. Comme les morphismes de transition dans le système inductif envisagé dans des immersions ouvertes, il est d'ailleurs clair que cette limite est en fait représentable par un schéma en groupes (en général non quasi-compact) sur $K$, que nous désignerons comme de \unsure{juste|juśte} par le même \unsure{symbole|symbôle} $A^{\handwritten{\flat}}_K$. On notera qu'on a, pour ce schéma en groupes
\[
A_K^{\handwritten{\flat}\,o}=A_K^{\handwritten{\natural}\,o}\ ,\
A_K^{\handwritten{\flat}}/A_K^{\handwritten{\flat}\,o}\ \simeq\
\underset{K'}{\mathop{\lim}\limits_{\handwritten{\longrightarrow}}}\Phi_{K,K'}\ .\tag*{(8.3.4)}
\]
Ce dernier groupe sera déterminé dans 11.10 : il est canoniquement isomorphe à $\underline{M}_K\otimes\mathbb Q/\mathbb Z$ (avec la notation (8.1.8)).

\handwritten{417}
\mspage{417}
% archive sheet number: 259
